\begin{figure}[H]
\centering
\begin{tikzpicture}[font=\small,>=Stealth]

\draw[->,gray] (-.5,0)--(8.6,0) node[right] {$x$};
\draw[->,gray] (0,-1.6)--(0,1.7) node[above] {$y$};
\draw[red!70!black,very thick] (0,-1.2)--(0,1.2);
\draw[blue,domain=.018:1,samples=2200] plot ({8*\x},{1.2*sin(180/pi/\x)});
\node[anchor=west] at (2,1.65) {$y=\sin(1/x)$};
\node[anchor=east] at (-.15,1.2) {$1$};\node[anchor=east] at (-.15,-1.2) {$-1$};
\node[below] at (8,0) {$1$};

\end{tikzpicture}
\caption{The vertical segment belongs to the closure of the oscillating graph. The drawing truncates the oscillations near zero; the proof uses all of them. A path leaving the segment would have to attain both heights arbitrarily close to one time.}
\label{fig:sine-curve}
\end{figure}
